\documentclass[11pt]{article}
\usepackage[T1]{fontenc}
\usepackage{lmodern}
\usepackage[margin=1in]{geometry}
\usepackage{amsmath,amssymb,amsthm,mathtools}
\usepackage{booktabs,array,longtable}
\usepackage{xcolor,microtype,fancyhdr}
\usepackage[hyphens]{url}
\usepackage[hypertexnames=false,colorlinks=true,linkcolor=blue!45!black,citecolor=blue!45!black,urlcolor=blue!45!black]{hyperref}
\usepackage{enumitem}
\usepackage{listings}
\setlist{topsep=4pt,itemsep=3pt}
\lstset{basicstyle=\ttfamily\footnotesize,breaklines=true,columns=fullflexible,keepspaces=true,frame=single,rulecolor=\color{black!15}}
\pagestyle{fancy}\fancyhf{}
\fancyhead[L]{\small Constructive calculus in iPRA}
\fancyhead[R]{\small Explicit finite proofs}
\fancyfoot[C]{\thepage}
\setlength{\headheight}{14pt}
\setlength{\emergencystretch}{2em}
\newtheorem{theorem}{Theorem}[section]
\newtheorem{lemma}[theorem]{Lemma}
\newtheorem{proposition}[theorem]{Proposition}
\newtheorem{corollary}[theorem]{Corollary}
\theoremstyle{definition}
\newtheorem{definition}[theorem]{Definition}
\newtheorem{remark}[theorem]{Remark}
\newcommand{\N}{\mathbb N}
\newcommand{\Z}{\mathbb Z}
\newcommand{\Q}{\mathbb Q}
\newcommand{\R}{\mathbb R}
\newcommand{\ipra}{\ensuremath{\mathsf{iPRA}}}
\newcommand{\eps}{\varepsilon}
\newcommand{\abs}[1]{\left|#1\right|}
\newcommand{\normc}[1]{\left\|#1\right\|_{\mathrm{coef}}}
\newcommand{\dd}{\,\mathrm{d}}
\newcommand{\mesh}{\operatorname{mesh}}
\newcommand{\sgn}{\operatorname{sgn}}
\newcommand{\PR}{\mathsf{PR}}
\newcommand{\code}[1]{\texttt{#1}}
\title{The fundamental theorem of calculus\protect\\and the Euler--Maclaurin formula\protect\\in intuitionistic primitive recursive arithmetic}
\author{A constructive proof package prepared for Vladimir Reshetnikov}
\date{November 14, 2026}
\hypersetup{pdftitle={Constructive calculus in intuitionistic primitive recursive arithmetic},pdfauthor={Prepared for Vladimir Reshetnikov},pdfsubject={Explicit finite proofs of FTC and Euler--Maclaurin, with a Lean formalization roadmap},pdfkeywords={iPRA, constructive calculus, fundamental theorem of calculus, Euler-Maclaurin, Bernoulli polynomials}}

\begin{document}
\maketitle
\begin{abstract}
We give complete human-readable proofs of an effective fundamental theorem
of calculus, in both directions, and the finite-order Euler--Maclaurin
formula.  The mathematical objects are primitive recursive rational
approximation procedures with explicit regularity certificates.  Every
analytic equality means an explicitly controlled family of rational
inequalities.  Integration is constructed from Riemann sums; it is not
defined by antiderivatives.  The fundamental theorem follows from finite
refinement estimates and telescoping.  Euler--Maclaurin follows from
rational Bernoulli polynomials and a finite product identity.  We give a
fully rational error estimate uniform in the truncation order, and the
elementary remainder bound
\(\abs{R_m}\le (b-a)4^{-m}M_{2m}\).

First-order arithmetic cannot quantify over arbitrary real functions.
Accordingly, the precise results are theorem schemas in fixed primitive
recursive evaluators, with stated moduli as hypotheses.  The usual
continuous-function form of FTC II uses a supplied uniform derivative
modulus.  A separate Riemann-integrable version uses a supplied primitive
recursive tagged-partition realizer.  These data are not extracted from
bare qualitative hypotheses by an unmentioned compactness principle.

The package includes exact-arithmetic implementations and checks, an
uncompiled Lean candidate for the finite arithmetic layer, and a detailed
route to object-theory proofs in ProveIt's deeply embedded \ipra.  It does
not claim a completed or kernel-checked Lean formalization.
\end{abstract}
\newpage
\tableofcontents
\newpage

\section{The results and their precise scope}\label{sec:scope}

The familiar formulas we prove are
\begin{align}
 A(x)&=\int_a^x f(t)\dd t \quad\Longrightarrow\quad A'(x)=f(x),
 \label{eq:ftc1intro}\\
 H'&=f \quad\Longrightarrow\quad
 \int_a^b f(t)\dd t=H(b)-H(a),\label{eq:ftc2intro}
\end{align}
and, for integers \(a\le b\) and \(m\ge1\),
\begin{align}
 \sum_{k=a}^{b}f(k)
 ={}&\int_a^b f(x)\dd x+\frac{f(a)+f(b)}2\notag\\
 &+\sum_{s=1}^{m}\frac{B_{2s}}{(2s)!}
       \bigl(f^{(2s-1)}(b)-f^{(2s-1)}(a)\bigr)+R_m,\label{eq:emintro}\\
 R_m={}&-\sum_{k=a}^{b-1}\int_0^1
       \frac{B_{2m}(t)}{(2m)!}f^{(2m)}(k+t)\dd t.\label{eq:rintro}
\end{align}
The remainder is part of an exact identity.  We make no assertion that
the sequence of corrections converges when \(m\) tends to infinity.
The convention is \(B_1=-1/2\).  Formula \eqref{eq:emintro} agrees with
the standard finite Euler--Maclaurin formula, for example the equivalent
version in NIST DLMF \S2.10(i) \cite{dlmfem}.

There is a necessary representational qualification.  The language of
\ipra\ has number variables, not variables for arbitrary functions
\(\R\to\R\), sets of real numbers, or arbitrary infinite sequences.
Here a function is supplied by a \emph{fixed primitive recursive}
procedure producing rational approximations, together with fixed
primitive recursive moduli and arithmetic correctness assumptions.
Each choice of those procedures gives a first-order arithmetic theorem.
Number parameters inside a fixed procedure are allowed uniformly.

This is substantive: the bounds below supply the number of subdivisions,
the precision of every evaluation, and the exact finite induction
invariants.  It would not establish the claimed logical strength merely
to prove an ordinary theorem about \(\R\) in Lean or to say that all the
objects are computable.

\begin{center}
\begin{tabular}{>{\raggedright\arraybackslash}p{.20\textwidth}>{\raggedright\arraybackslash}p{.70\textwidth}}
\toprule
Result & Explicit input data\\\midrule
FTC I & A rational approximation procedure, a uniform continuity modulus,
and a bound on a rational compact interval.\\
FTC II, uniform form & The preceding data and an approximation procedure
for an antiderivative with a uniform first-order remainder modulus.\\
FTC II, Riemann form & A tagged-sum Cauchy modulus, pointwise derivative
certificates, and a primitive recursive procedure producing suitable
finite tagged partitions.\\
Euler--Maclaurin & A finite derivative jet, or a fixed primitive recursive
derivative family, with bounds, continuity moduli, and the adjacent
derivative certificates.\\
\bottomrule
\end{tabular}
\end{center}

The uniform form is the effective \(C^1\) version used in the main
construction.  It is not presented as a proof that a qualitative
pointwise derivative assertion automatically has a primitive recursive
uniform modulus.  Section~\ref{sec:generalFTC} states the additional
partition data needed for the more general Riemann-integrable version.
In particular, we do not claim the unrestricted set-theoretic theorems
as single sentences of first-order \ipra.

\section{The arithmetic foundation}\label{sec:foundation}

\subsection{Logic, induction, and coding}

Use first-order intuitionistic logic with equality, \(0\), successor,
and a symbol for each primitive recursive function.  Include their
defining equations, successor injectivity, \(Sx\ne0\), and induction for
quantifier-free formulas, with parameters.  This is the ordinary
presentation of \ipra\ formalized by \code{PRAAx} in ProveIt's
\code{ConciseIPRA/PR.lean} \cite{prsource}.  Its underlying proof system
is intuitionistic natural deduction.  No unrestricted excluded middle,
choice principle, Markov principle, or induction for arbitrary formulas
is used here.

Signed integers are coded by natural numbers, for example by a sign and
magnitude.  A rational is a pair consisting of a signed numerator and a
positive denominator; store the denominator as \(d+1\).  Rational
comparison is signed-integer cross multiplication.  Intermediate
arithmetic pairs need not be reduced, but every rational argument passed
to an input evaluator is first put in a fixed canonical representation.
Reduction uses the greatest common divisor, obtainable by bounded search
through the positive common divisors, and is primitive recursive.
Thus equal rational grid nodes give literally the same evaluator input,
including shared cell endpoints.  Finite lists can be coded by primitive
recursive pairing and projection functions, with an explicit length.
All lists, polynomials, partitions, sums, and bounded maxima in this
paper are such finite codes.  Rational division is used only when the
denominator has an explicitly verified positive or nonzero rational
value.

Every bounded quantifier in a proposed induction invariant can be
replaced by a primitive recursive Boolean fold.  Thus
\(\forall i<n\,C(i)\), for a decidable arithmetic test \(C\), becomes
\(\operatorname{all}_{i<n}\chi_C(i)=1\).  A finite triangle inequality,
a polynomial coefficient identity, or a bound on a finite list then
has a genuinely quantifier-free formulation.  The universal data
assumptions stay in the side context and are instantiated at the
indices occurring in each step.  They are not the induction formula.

Write
\[
 \eps_k=2^{-k}.
\]
For a nonnegative rational \(z\), let \(\sigma(z)\) be a natural number
such that
\begin{equation}
 2^{\sigma(z)}\ge\max(1,z).\label{eq:sigma}
\end{equation}
For definiteness take the least such number within
\(0,\ldots,\lceil z\rceil+1\).  This is bounded search, hence primitive
recursive; the bound exists since \(2^n\ge n+1\), by quantifier-free
induction.  No logarithm oracle is being assumed.  The notation
\(\sigma\) merely makes precision schedules readable.

\subsection{A schema, not a universal evaluator}

Expressions such as \(F(x,p)\) and \(\mu(k)\) stand for fixed primitive
recursive definitions.  A displayed proof is a proof template under
substitution of those definitions.  We do not add a supposedly primitive
recursive evaluator for all primitive recursive programs.  Such a
universal evaluator is not available.  Relative primitive recursion is
only metanotation here: substituting fixed primitive recursive input
maps into the constructions produces ordinary primitive recursive
maps.

For a fixed order \(m\), the Euler--Maclaurin schema can take finitely
many independently supplied derivative evaluators.  If \(m\) is itself
a number variable, use one fixed map \(F(j,x,p)\), and fixed maps for
its bounds and moduli, with derivative index \(j\) as an additional
number argument.  The assumptions for a chosen order concern only
\(j\le2m\).  Section~\ref{sec:finiteEM} proves the required uniform
finite estimate without induction on an analytic assertion.

\subsection{Eliminating limiting shorthand}\label{sec:closure}

\begin{lemma}[Rational closure]
If rationals \(u,v,C\), with \(C\ge0\), satisfy
\(u\le v+C2^{-n}\) for every natural number \(n\), then \(u\le v\).
\end{lemma}
\begin{proof}
Compare the rationals \(u,v\), a decidable operation.  If \(u>v\), write
\(u-v=w/d\) with integers \(w,d\ge1\), and choose an integer
\(D\ge\max(1,C)\).  At the explicitly bounded index \(n=Dd+1\),
\(2^n\ge n+1>Dd\), so \(C2^{-n}<1/d\le u-v\), a contradiction.
The other branch is already \(u\le v\).  This is rational decidability,
not excluded middle for order on real names.
\end{proof}

When a proof below says that an auxiliary error tends to zero, first
fix the output approximation index.  The remaining assertion is an
inequality between rationals, and the lemma eliminates the auxiliary
error.  Alternatively, select that error below any requested dyadic
tolerance using \eqref{eq:sigma}.  No completeness axiom or unbounded
search is needed.  When an equality of represented quantities is used
inside a finite sum, the indices can be chosen simultaneously by a
bounded maximum.  Section~\ref{sec:finiteEM} gives this selection in full
for the longest argument in the paper.

\section{Names and certified function data}\label{sec:names}

\begin{definition}[Regular rational name]
A regular name is a fixed primitive recursive sequence \(X(p)\in\Q\)
with the arithmetic promise
\begin{equation}
 \abs{X(p)-X(q)}\le\eps_p+\eps_q\qquad(p,q\in\N).\label{eq:name}
\end{equation}
Two names represent the same number if
\(\abs{X(p)-Y(q)}\le\eps_p+\eps_q\) for all \(p,q\).
\end{definition}

Reflexivity and symmetry are immediate; transitivity follows by putting
an intermediate name at precision \(n\), applying the triangle
inequality, and using rational closure.  A rational is represented by
its constant sequence.  Addition or subtraction is represented by
\(X(p+1)\pm Y(p+1)\).  If bounds \(\abs X\le M_X\) and
\(\abs Y\le M_Y\) are supplied, multiplication is represented by
\(X(s)Y(s)\), where
\(s=p+\sigma(2(M_X+M_Y+1))\).  Indeed, its evaluation error is bounded
by \((M_X+M_Y)\eps_s+\eps_s^2\le\eps_p/2\).
These statements can all be read as cross-index rational inequalities.

Notation such as \(\abs{X-c}\le d\), for rational \(c,d\), means
\(\abs{X(p)-c}\le d+\eps_p\) for every \(p\).
Nonnegativity \(X\ge0\) means \(X(p)\ge-\eps_p\) for every \(p\).
Likewise \(\abs{X-Y}\le d\) means the corresponding cross-index
inequality with \(\eps_p+\eps_q\) added.  Real-number notation below
is an abbreviation for this arithmetic, not an additional object sort.

Fix a rational interval \([a,b]\), with \(a\le b\), and set \(L=b-a\).
Choose once and for all \(s=\sigma(L)\).  The same \(s\) is used for
every subinterval.  An admissible uniformly continuous function consists
of a fixed primitive recursive rational evaluator \(F(x,p)\), a
primitive recursive modulus \(\mu\), and an integer \(M\ge1\), with
the following promises for rational arguments in \([a,b]\):
\begin{align}
 \abs{F(x,p)-F(x,q)}&\le\eps_p+\eps_q,\tag{C}\label{eq:C}\\
 \abs{F(x,p)}&\le M+\eps_p,\tag{B}\label{eq:B}\\
 \abs{x-y}\le2^{-\mu(k)}\quad&\Longrightarrow\quad
 \abs{F(x,p)-F(y,q)}\le\eps_k+\eps_p+\eps_q.\tag{U}\label{eq:U}
\end{align}
We write \(f(x)\) for the name \(p\mapsto F(x,p)\).  Normalize the
modulus, if necessary, by
\begin{equation}
 \mu(k)=k+\max_{i\le k}\mu_{\mathrm{old}}(i).\label{eq:munormal}
\end{equation}
This preserves (U), is primitive recursive, and satisfies
\(\mu(k+1)\ge\mu(k)+1\) and \(\mu(k)\ge k\).

The bound (B) can be derived from (U), so it does not hide a compactness
theorem.  Let \(N=\max(1,\lceil L2^{\mu(0)}\rceil)\).  Subdivide
\([a,x]\) into \(N\) equal pieces.  At common precision \(p\), each
increment of \(F\) has magnitude at most \(1+2\eps_p\le3\), and
\(\abs{F(a,p)}\le\abs{F(a,0)}+1+\eps_p\).  Thus
\(M=\max(1,\lceil\abs{F(a,0)}\rceil+1+3N)\) is a valid explicit bound.

\begin{definition}[Uniform derivative certificate]\label{def:derivative}
Let \(H(x,p)\) be another coherent evaluator.  A primitive recursive
\(\nu\) certifies that \(H'=f\) if, whenever
\(\abs{y-x}\le2^{-\nu(k)}\),
\begin{equation}
 \abs{H(y,p)-H(x,p)-(y-x)F(x,p)}
 \le\abs{y-x}\eps_k+(2+\abs{y-x})\eps_p.\tag{D}\label{eq:D}
\end{equation}
All arguments in this promise are rational and belong to the interval.
\end{definition}

The last term is essential: two antiderivative evaluations and one
derivative evaluation are being rounded.  In name notation (D) says
\[
 \abs{H(y)-H(x)-(y-x)f(x)}\le\abs{y-x}\eps_k.
\]
It is a division-free quantitative definition of differentiability.
For a nonzero rational increment it implies the usual difference-quotient
estimate.  The case \(x=y\) needs no special decision on real equality.

\section{Constructing the integral from finite sums}\label{sec:integral}

For rationals \(u\le v\) in the interval and an integer \(N\ge1\), set
\begin{equation}
 S_{N,p}(u,v)=\frac{v-u}{N}\sum_{i=0}^{N-1}
 F\!\left(u+i\frac{v-u}{N},p\right).\label{eq:sum}
\end{equation}
This is a primitive recursive rational expression.

\subsection{The finite refinement estimate}

If a partition is refined and each old sample is repeated over its
subcells, the weighted sum is unchanged: the subcell lengths sum to the
parent length.  This is a finite distributive identity.  If instead
each repeated sample is replaced by a sample at a nearby tag, (U) bounds
the change by the total length times the common pointwise error.

For two uniform left sums, use their common subdivision into \(NM\)
pieces.  On a subcell, the two relevant left endpoints both lie to the
left of its points; their distance is at most the larger original mesh.
Consequently, if both meshes are at most \(2^{-\mu(k)}\),
\begin{equation}
 \abs{S_{N,p}(u,v)-S_{M,q}(u,v)}
 \le (v-u)(\eps_k+\eps_p+\eps_q).\label{eq:refine}
\end{equation}
The same argument works for arbitrary rational partitions with left
tags, by their finite common refinement.  Only comparisons of rational
endpoints are needed.

For completeness, a refinement can be avoided altogether.  For cells
\([x_i,x_{i+1}]\) and \([y_j,y_{j+1}]\), assign the rational overlap
weight
\[
 w_{ij}=\max\bigl(0,\min(x_{i+1},y_{j+1})-\max(x_i,y_j)\bigr).
\]
Its row and column sums are the respective cell lengths.  For a fixed
row interval \([\alpha,\beta]\), put
\(C(t)=\min(\beta,\max(\alpha,t))\).  Rational order cases give
\(w_{ij}=C(y_{j+1})-C(y_j)\), so the row sum telescopes to
\(\beta-\alpha\); the column proof is identical.  These are
quantifier-free finite inductions.  Rewrite both sums with these weights
and subtract.  For arbitrary tags, overlapping cells give tag distance
at most the sum of the two meshes.  Thus meshes at most
\(2^{-\mu(k+1)}\) suffice for the same estimate, using
\eqref{eq:munormal}.  This explicit observation will also justify
additivity and tagged-sum comparisons.

\subsection{A canonical primitive recursive name}

Put
\begin{equation}
 c_k=k+s+4,\qquad N_k=2^{s+\mu(c_k)},\qquad
 I_{u,v}(k)=S_{N_k,c_k}(u,v).\label{eq:integralname}
\end{equation}
The mesh is at most \(2^{-\mu(c_k)}\).  If \(k\le l\), \(N_l\) is
a multiple of \(N_k\), and \eqref{eq:refine} gives
\begin{align}
 \abs{I_{u,v}(k)-I_{u,v}(l)}
 &\le(v-u)(2\eps_{c_k}+\eps_{c_l})\notag\\
 &\le\frac{\eps_k}{8}+\frac{\eps_l}{16}.\label{eq:integralcauchy}
\end{align}
For \(l<k\), interchange them.  In particular this is a regular name.
Define \(\int_u^v f\) to be this name, \(\int_u^u f=0\), and
\(\int_v^u f=-\int_u^v f\).  No existence theorem for limits is used:
\eqref{eq:integralname} is the definition of the output procedure, and
\eqref{eq:integralcauchy} is its proved arithmetic certificate.

If \((v-u)/N\le2^{-\mu(k)}\), comparison with sufficiently fine
canonical sums and rational closure yield
\begin{equation}
 \abs{\int_u^v f-S_{N,p}(u,v)}
 \le(v-u)(\eps_k+\eps_p).\label{eq:quaderror}
\end{equation}
To read this entirely arithmetically, fix an index \(j\) of the name
on the left and add \(\eps_j\) to the right.  Apply
\eqref{eq:refine} with a canonical sum at index \(l\), use
\eqref{eq:name} between indices \(j,l\), and eliminate its explicitly
vanishing errors by Section~\ref{sec:closure}.  This proves every
rational instance of \eqref{eq:quaderror}.

\subsection{Basic integral rules}

\begin{proposition}\label{prop:integralrules}
The constructed integral is additive in adjacent rational intervals,
linear over the rationals, and satisfies
\[
 \abs{\int_u^v f}\le M\abs{v-u}.
\]
If \(f\ge0\), then \(\int_u^v f\ge0\) for \(u\le v\).
If \(\abs{f-g}\le C\), then
\(\abs{\int_u^v f-\int_u^v g}\le C\abs{v-u}\).
\end{proposition}
\begin{proof}
For additivity, compare the canonical partition of \([u,v]\) at index
\(k\) with the union of the canonical partitions of \([u,w]\) and
\([w,v]\), using the same global \(s\).  All sample precisions are
\(c_k\).  The overlap argument with arbitrary tags, at precision
\(c_k-1\), gives
\begin{equation}
 \abs{I_{u,v}(k)-I_{u,w}(k)-I_{w,v}(k)}
 \le4(v-u)\eps_{c_k}\le\eps_k/4.\label{eq:addfinite}
\end{equation}
This all-precision estimate proves equality of the represented names,
after the fixed shift used for adding names.  Reversing endpoints gives
all other orderings.

Linearity is the distributive law for a common finite partition.  Use
\eqref{eq:quaderror} to compare that partition with each canonical
integral, adjusting sample precisions for scalar factors by
\eqref{eq:sigma}.  The finite sums for a bounded or nonnegative
function have the asserted bounds, with the evaluation errors added.
Rational closure removes those errors.  The last assertion follows by
linearity and the same estimate applied to \(f-g\).
\end{proof}

The rules also apply to products of bounded uniformly continuous data.
For example, if \(f,g\) have bounds \(M_f,M_g\), the elementary identity
\[
 f(x)g(x)-f(y)g(y)
 =f(x)(g(x)-g(y))+g(y)(f(x)-f(y))
\]
provides a product modulus
\[
 \mu_{fg}(k)=\max\{\mu_f(k+\sigma(2M_g)),
                       \mu_g(k+\sigma(2M_f))\}.
\]
Evaluating both factors with the multiplication shift from
Section~\ref{sec:names} provides the required name.  Sums and scalar
products have equally explicit moduli.  These facts establish the
integrability of every product used later.

\section{The fundamental theorem of calculus}\label{sec:FTC}

\subsection{FTC I: integration produces a certified derivative}

\begin{theorem}[Effective FTC I]\label{thm:ftc1}
From (C), (B), and (U), the primitive recursive evaluator
\[
 G(x,p)=I_{a,x}(p)
\]
defines \(A(x)=\int_a^x f\), and satisfies the derivative certificate
\eqref{eq:D} with \(H=G\) and \(\nu=\mu\).
For \(a<b\), this means \(A'=f\), with one-sided derivatives at the
interval endpoints.  For \(a=b\), the zero integral and division-free
certificate remain valid, without asserting an ordinary derivative on
a singleton domain.
\end{theorem}
\begin{proof}
Let \(h=y-x\) and \(\abs h\le2^{-\mu(k)}\).  Every point between
\(x,y\) is within that distance of \(x\).  For each summand of the
canonical oriented integral over this interval, (U) gives
\[
 \abs{F(t,c_p)-F(x,p)}\le\eps_k+\eps_{c_p}+\eps_p.
\]
Multiplying by the cell lengths and adding yields the exact rational
bound
\begin{equation}
 \abs{I_{x,y}(p)-hF(x,p)}
 \le\abs h(\eps_k+\eps_{c_p}+\eps_p).\label{eq:localintegral}
\end{equation}
Here \(I_{x,y}\) is signed if \(y<x\).  By
\eqref{eq:addfinite},
\[
 \abs{G(y,p)-G(x,p)-I_{x,y}(p)}\le\eps_p/4.
\]
Also \(\abs h\le1\), since \(\mu(k)\ge0\), and
\(\eps_{c_p}\le\eps_p/16\).  Combining these facts gives
\[
 \abs{G(y,p)-G(x,p)-hF(x,p)}
 \le\abs h\eps_k+(2+\abs h)\eps_p,
\]
as asserted.  This argument already proves the literal rational
certificate; it does not first appeal to a theorem about real limits.

In name notation the same argument is the familiar estimate
\begin{equation}
 \abs{A(y)-A(x)-f(x)(y-x)}
 =\abs{\int_x^y(f(t)-f(x))\dd t}
 \le\abs{y-x}\eps_k.\label{eq:ftc1error}
\end{equation}
At nonzero rational \(h\), division by \(\abs h\) gives the
difference-quotient bound.  At \(h=0\), the division-free inequality
remains valid.  No integral mean-value theorem has been used.
\end{proof}

Proposition~\ref{prop:integralrules} also gives
\(\abs{A(y)-A(x)}\le M\abs{y-x}\), so the constructed primitive is
Lipschitz.  Its continuity modulus can be taken as \(k+\sigma(M)\)
in name notation, with a harmless extra shift for the evaluator.

\subsection{FTC II: local errors telescope}

\begin{theorem}[Effective uniform FTC II]\label{thm:ftc2}
Suppose \(f\) has (C), (B), (U), and a coherent evaluator \(H\) has
the supplied derivative certificate (D).  Then
\[
 \int_a^b f=H(b)-H(a)
\]
as equality of names.  Every requested precision has an explicit
primitive recursive finite proof.
\end{theorem}
\begin{proof}
If \(L=0\), both sides are zero.  Suppose \(L>0\).  On a uniform
partition \(x_i=a+ih\), \(h=L/N\), assume
\(h\le2^{-\nu(k)}\).  At common evaluation precision \(p\), sum (D)
over the \(N\) cells.  The exact same samples at interior endpoints
cancel, so
\begin{equation}
 \abs{H(b,p)-H(a,p)-S_{N,p}(a,b)}
 \le L\eps_k+(2N+L)\eps_p.\label{eq:finiteFTC}
\end{equation}
This is a finite triangle inequality followed by the algebraic identity
\(\sum_{i<N}(H(x_{i+1},p)-H(x_i,p))=H(b,p)-H(a,p)\).
The term \(2N\eps_p\) records the local evaluation errors; it must
not be dropped just because endpoint samples telescope.

For a requested accuracy \(\eps_r\), choose
\begin{align}
 k&=r+s+4,\notag\\
 N&=2^{s+\max(\mu(k),\nu(k))},\notag\\
 p&=r+\sigma(2N+2L+2)+4.\label{eq:FTCprecision}
\end{align}
The mesh meets both moduli.  Combine \eqref{eq:finiteFTC},
\eqref{eq:quaderror}, and the two endpoint evaluation errors.  The
resulting bound on the represented difference is
\begin{align*}
 \abs{H(b)-H(a)-\int_a^b f}
 &\le 2L\eps_k+(2N+2L+2)\eps_p\\
 &\le\eps_r/8+\eps_r/16<\eps_r.
\end{align*}
Every term in \eqref{eq:FTCprecision} is primitive recursive.  Since
\(r\) was arbitrary, the names are equal.  More literally, fix the
indices of the names being compared, retain their radius terms, and
use rational closure on the auxiliary \(r\).  This proves the
cross-index definition of equality from Section~\ref{sec:names}.
\end{proof}

\subsection{Arguments given as real names}

The rational-point statements determine the usual effective real
functions.  Here are the explicit details.  Let \(X(n)\) be a fixed
primitive recursive regular name for an interval point, with its
rational approximants clamped to \([a,b]\).  Rational clamping is
primitive recursive and does not change a name representing an interval
point.  Define
\begin{equation}
 \widehat F(X,p)=F\bigl(X(\mu(p+2)+2),p+2\bigr).\label{eq:extend}
\end{equation}
For \(p\le q\), the two selected rational arguments differ by at most
\(2\eps_{\mu(p+2)+2}\le2^{-\mu(p+2)}\).  Hence (U) gives
\[
 \abs{\widehat F(X,p)-\widehat F(X,q)}
 \le2\eps_{p+2}+\eps_{q+2}.
\]
This is a regular name.  The same calculation for equivalent input
names proves that \eqref{eq:extend} respects represented equality;
the output discrepancy tends to zero with an explicit index shift.
The Lipschitz primitive \(A\) is extended by the identical construction
using its computed modulus.

To transfer \eqref{eq:ftc1error} to two named interval points, leave a
dyadic margin.  The modulus
\(\mu(k+1)+1\) ensures that sufficiently accurate rational
representatives are within \(2^{-\mu(k+1)}\), under the antecedent
\(\abs{Y-X}\le2^{-(\mu(k+1)+1)}\).  Apply the rational
estimate and let only the explicitly bounded approximation errors
decrease.  Rational closure, at fixed output indices, yields
\[
 \abs{A(Y)-A(X)-f(X)(Y-X)}\le\eps_k\abs{Y-X}.
\]
Thus FTC I holds at every admissible primitive recursive named argument.
Difference quotients for a named nonzero increment require an
\emph{apartness witness}, such as \(\abs h\ge2^{-d}>0\), to compute
its reciprocal.  The division-free statement requires no such
witness, and never decides real trichotomy.  This extension remains a
schema in fixed names, not an internal quantification over all real
sequences.

\section{The Riemann-integrable form of FTC II}\label{sec:generalFTC}

The usual weakest elementary version of FTC II assumes that the
derivative is Riemann integrable, rather than continuous.  The uniform
modulus in Definition~\ref{def:derivative} is stronger data than bare
pointwise differentiability.  We therefore state a separate precise
arithmetic theorem that exposes the partition information needed in
that version.

The case \(L=0\) is the identity \(0=H(a)-H(a)\), so assume \(L>0\)
in this section.  Let \(F\) satisfy coherence and boundedness on
rational arguments.  A
tagged rational partition \(P\) consists of
\(a=x_0<\cdots<x_N=b\), with rational tags
\(t_i\in[x_i,x_{i+1}]\).  Set
\[
 S(P,p)=\sum_{i<N}(x_{i+1}-x_i)F(t_i,p).
\]
Supply a primitive recursive tagged-sum modulus \(\rho\) satisfying
\begin{equation}
 \begin{split}
 \mesh(P),\mesh(Q)\le2^{-\rho(k)}\quad\Longrightarrow\quad
 \abs{S(P,p)-S(Q,q)}\le\eps_k+L(\eps_p+\eps_q).
 \end{split}\tag{R}\label{eq:Riemann}
\end{equation}
The quantifiers over \(P,Q\) are quantifiers over finite natural-number
codes with validity tests.  Replace \(\rho(k)\), if necessary, by
\(\max_{i\le k}\rho_{\mathrm{old}}(i)\), so it is nondecreasing;
this preserves (R).  Canonical uniform sums, with mesh controlled
by \(\rho(k+4)\) and sample precision \(k+s+4\), give a regular
primitive recursive integral name.  Indeed for indices \(k\le l\),
\eqref{eq:Riemann} bounds their difference by
\(\eps_{k+4}+L(\eps_{k+s+4}+\eps_{l+s+4})\), which is smaller than
\(\eps_k+\eps_l\).  Any tagged sum with mesh at most
\(2^{-\rho(k)}\) approximates this integral within
\(\eps_k+L\eps_p\), by rational closure.

This is a rational evaluation interface.  If one starts externally with
an ordinary Riemann-integrable real function and also supplies a
primitive recursive evaluator and a primitive recursive tagged-sum
modulus with the stated certificates, its modulus supplies (R) on
rational tags and the represented integral is its usual integral.
Qualitative Riemann integrability alone does not supply those primitive
recursive data.  We do not infer ordinary Riemann integrability of
an arbitrary discontinuous real extension merely from its values on
the rationals.  Unlike the uniform-continuity case, those values alone
need not determine a real function.

Suppose the antiderivative evaluator \(H\) is coherent, in the sense of
(C), and has the pointwise certificate
\begin{equation}
 \begin{split}
 \abs{u-t}\le2^{-\nu(t,k)}\quad\Longrightarrow\quad
 \abs{H(u,p)-H(t,p)-(u-t)F(t,p)}\\
 \le\abs{u-t}\eps_k+(2+\abs{u-t})\eps_p
 \end{split}\tag{PD}\label{eq:PD}
\end{equation}
at rational tags.  Finally, supply a primitive recursive procedure
\(\mathcal P(k)\) returning a valid tagged partition with
\begin{equation}
 \mesh(\mathcal P(k))\le2^{-\rho(k)},\qquad
 \abs{x_i-t_i},\abs{x_{i+1}-t_i}\le2^{-\nu(t_i,k)}.
 \tag{P}\label{eq:partitionwitness}
\end{equation}
The existence and correctness of this procedure are explicit hypotheses.

\begin{theorem}[FTC II with effective tagged divisions]\label{thm:taggedFTC}
Under (R), (PD), and (P),
\(\int_a^b f=H(b)-H(a)\) as equality of names.
\end{theorem}
\begin{proof}
Apply (PD) at tag \(t_i\) to both endpoints of its cell and subtract.
The two values \(H(t_i,p)\) cancel exactly.  Since the tag lies in the
cell,
\(\abs{x_i-t_i}+\abs{x_{i+1}-t_i}=x_{i+1}-x_i\).
Thus the cell error is at most
\((x_{i+1}-x_i)\eps_k+(4+x_{i+1}-x_i)\eps_p\).
Summing and telescoping gives
\begin{equation}
 \abs{H(b,p)-H(a,p)-S(\mathcal P(k),p)}
 \le L\eps_k+(4N+L)\eps_p.\label{eq:taggedfinite}
\end{equation}
Using (R) and the two endpoint radii, the represented error is at most
\[
 (L+1)\eps_k+(4N+2L+2)\eps_p.
\]
For target precision \(r\), take
\[
 k=r+\sigma(L+1)+4,\qquad
 p=r+\sigma(4N+2L+2)+4,
\]
where \(N\) is the length of the already computed
\(\mathcal P(k)\).  Each term is at most \(\eps_r/16\).
All computation is primitive recursive, and all reasoning is finite.
\end{proof}

Under a uniform derivative modulus, (P) is supplied by a sufficiently
fine uniform grid, so this theorem includes the previous situation.
For a more general function an adaptive partition algorithm, with its
proved primitive recursive bound, may supply (P).

One must not discard (P) and cite a qualitative Cousin lemma or compactness
principle.  Even an arbitrary positive gauge given only on rational
tags need not have a rational tagged fine partition: externally, choose
each gauge radius smaller than the distance of its rational tag from a
fixed irrational point.  No cell tagged within such a radius can contain
that point.  The finite certificate is therefore a material condition,
not an automatic consequence of the rational pointwise statement.
This explicitly describes both the strength and the limitation of the
general Riemann version proved here.

\section{Integration by parts from finite algebra}\label{sec:parts}

For rational arrays \(u_i,v_i\), put
\(\Delta u_i=u_{i+1}-u_i\) and \(\Delta v_i=v_{i+1}-v_i\).
The exact identity
\begin{equation}
 u_{i+1}v_{i+1}-u_iv_i
 =u_i\Delta v_i+v_i\Delta u_i+\Delta u_i\Delta v_i
 \label{eq:productfinite}
\end{equation}
is just multiplication and subtraction.  Its sum telescopes.  This
identity is the entire algebra behind integration by parts.

Here is a quantitative proof that it gives the usual product rule in
our framework.  Suppose \(U,V\) are admissible uniformly continuous
function data, and they and their certified derivatives
\(u,v\) have respective bounds \(M_U,M_V,A,B\), each at least one.
For \(h=y-x\), write
\[
 E_U=U(y)-U(x)-hu(x),\qquad
 E_V=V(y)-V(x)-hv(x).
\]
If the derivative moduli are met at precision \(t\),
\(\abs{E_U},\abs{E_V}\le\abs h\eps_t\), and
\(\abs{\Delta U}\le\abs h(A+1)\),
\(\abs{\Delta V}\le\abs h(B+1)\).  Subtract the proposed linear
part from \eqref{eq:productfinite}:
\[
 \Delta(UV)-h(Uv+Vu)=U(x)E_V+V(x)E_U+\Delta U\Delta V.
\]
Choose
\begin{align*}
 t&=k+\sigma(4(M_U+M_V+1)),\\
 \lambda(k)&=\max\{\nu_U(t),\nu_V(t),
                   k+\sigma(4(A+1)(B+1))\}.
\end{align*}
For \(\abs h\le2^{-\lambda(k)}\) the magnitude of the last display
is at most
\[
 \abs h\bigl((M_U+M_V)\eps_t+(A+1)(B+1)\abs h\bigr)
 \le\tfrac12\abs h\eps_k.
\]
Name evaluation introduces precisely the radius terms of (D), with
the multiplication shift already specified.  Thus \(UV\) has the
certified derivative \(Uv+Vu\).  The derivative is integrable by the
product and sum moduli.  Applying FTC II and linearity gives
\begin{equation}
\int_\alpha^\beta U(t)v(t)\dd t
 =U(\beta)V(\beta)-U(\alpha)V(\alpha)
       -\int_\alpha^\beta u(t)V(t)\dd t.\label{eq:IBP}
\end{equation}

The argument used only fixed-precision instances of
\eqref{eq:productfinite} and the already proved FTC.  In
Section~\ref{sec:finiteEM} we supply a direct rational version for all
the products in Euler--Maclaurin at once; this removes any concern about
repeating name-level equalities under induction on the order.

\section{Bernoulli polynomials by rational recursion}\label{sec:bernoulli}

It is convenient to normalize away the factorials:
\[
 P_n(t)=\frac{B_n(t)}{n!},\qquad c_n=P_n(0)=\frac{B_n}{n!}.
\]
We construct the polynomials without exponential functions, formal
infinite series, or analytic integration.  Define
\begin{equation}
 P_0=1,\qquad
 A_n(t)=\int_0^t P_{n-1}(u)\dd u,\qquad
 P_n(t)=A_n(t)-\int_0^1 A_n(u)\dd u.\label{eq:Pdef}
\end{equation}
Both integrals in this definition mean \emph{coefficientwise rational
polynomial operations}.  For \(\sum_i a_it^i\), the first sends
\(a_i\) to \(a_i/(i+1)\) in degree \(i+1\); evaluation of that
antiderivative at rational endpoints is finite arithmetic.  In
particular, if \(P_{n-1}(t)=\sum_{i=0}^{n-1}a_it^i\), then
\begin{equation}
 [t^{i+1}]P_n=\frac{a_i}{i+1},\qquad
 [t^0]P_n=-\sum_{i=0}^{n-1}\frac{a_i}{(i+1)(i+2)}.\label{eq:Pcoeff}
\end{equation}
This is primitive recursion on a finite coefficient table.

\subsection{The required identities, with proofs}

\begin{lemma}\label{lem:Pproperties}
For \(n\ge1\),
\begin{align}
 P_n'&=P_{n-1},&\int_0^1P_n(t)\dd t&=0,\label{eq:Pderiv}\\
 P_n(1-t)&=(-1)^nP_n(t).&&\label{eq:Preflection}
\end{align}
Moreover \(P_1(t)=t-1/2\), and for \(n\ge2\),
\begin{equation}
 P_n(1)=P_n(0).\label{eq:Pendpoint}
\end{equation}
For odd \(n\ge3\), \(c_n=0\).
\end{lemma}
\begin{proof}
The first two assertions and the formula for \(P_1\) follow immediately
by the coefficient operations in \eqref{eq:Pdef}.  A polynomial with a
specified derivative and specified mean over \([0,1]\) is unique:
the difference has all positive-degree coefficients zero, so is
constant, and its mean makes that constant zero.

For reflection, use induction on the finite coefficient identity.
It holds for \(P_0=1\).  If it holds at \(n-1\), the polynomial
\(Q(t)=(-1)^nP_n(1-t)\) has derivative
\[
 Q'(t)=(-1)^{n+1}P_{n-1}(1-t)=P_{n-1}(t).
\]
Its coefficientwise integral is zero: the change of variable
\(u=1-t\) for a polynomial is a finite binomial identity (or follows
by differentiating its explicit antiderivative at the endpoints).
Uniqueness proves \(Q=P_n\), hence \eqref{eq:Preflection}.

For \(n\ge2\), coefficientwise integration of the derivative gives
\(P_n(1)-P_n(0)=\int_0^1P_{n-1}=0\).  If \(n\) is odd and at least
three, reflection gives \(P_n(1)=-P_n(0)\), so \(2P_n(0)=0\).
Rational arithmetic yields \(c_n=0\).
\end{proof}

To identify the convention with the usual Bernoulli numbers, repeated
use of \(P_n'=P_{n-1}\) gives the finite identity
\[
 P_n(t)=\sum_{j=0}^n c_j\frac{t^{n-j}}{(n-j)!},
 \qquad c_0=1.
\]
Its zero mean implies
\[
 c_n=-\sum_{j=0}^{n-1}\frac{c_j}{(n-j+1)!}.
\]
Multiplying by \((n+1)!\) gives
\(\sum_{j=0}^{n}\binom{n+1}{j}B_j=0\) for \(n\ge1\), the
standard triangular Bernoulli recurrence.  Thus these are the standard
polynomials and numbers, with \(B_1=-1/2\), consistently with
\cite{dlmfbernoulli,dlmfproperties}.

For reference,
\begin{align*}
 P_1(t)&=t-\tfrac12,\\
 P_2(t)&=\tfrac12t^2-\tfrac12t+\tfrac1{12},\\
 P_3(t)&=\tfrac16t^3-\tfrac14t^2+\tfrac1{12}t,\\
 P_4(t)&=\tfrac1{24}t^4-\tfrac1{12}t^3+\tfrac1{24}t^2-\tfrac1{720}.
\end{align*}
In particular \(c_2=1/12\), \(c_4=-1/720\), and \(c_6=1/30240\).

\subsection{Polynomial bounds and derivative certificates}

For \(P=\sum a_it^i\), put \(\normc P=\sum_i\abs{a_i}\).
For rational \(t\in[0,1]\),
\(\abs{P(t)}\le\normc P\).  From \eqref{eq:Pcoeff},
\begin{equation}
 \normc{P_n}\le\tfrac32\normc{P_{n-1}}
 \le(3/2)^n\le2^n.\label{eq:coefbound}
\end{equation}
The first inequality follows because the nonconstant coefficients
contribute at most \(\sum\abs{a_i}\), and the new constant coefficient
at most \(\tfrac12\sum\abs{a_i}\).

No analytic mean-value theorem is needed for the polynomial derivative
bounds.  For \(0\le t\le t+h\le1\), factor
\[
 (t+h)^d-t^d=h\sum_{v=0}^{d-1}(t+h)^{d-1-v}t^v.
\]
Every factor in the sum has magnitude at most one.  Subtracting
\(dh t^{d-1}\), factoring again, and summing
\(0+1+\cdots+(d-1)=d(d-1)/2\), gives
\begin{equation}
 \abs{(t+h)^d-t^d-dht^{d-1}}
 \le\tfrac12d(d-1)h^2.\label{eq:monomialremainder}
\end{equation}
The same bound with \(\abs h\) holds for reversed endpoints, by
the algebraic factorization with \(x,y\in[0,1]\).  Summing by
coefficients proves a Lipschitz bound and a uniform derivative
certificate for every rational polynomial.  Thus the coefficientwise
antiderivatives used above also have the corresponding analytic
integrals by FTC II; the two meanings of polynomial integration agree.

\section{Euler--Maclaurin: the human-readable proof}\label{sec:EM}

Let \(a,b\in\Z\), \(a\le b\), \(r\ge1\), and let
\(f_0,f_1,\ldots,f_r\) be admissible function data on \([a,b]\), with
\(f_{j-1}'=f_j\) certified by (D), for \(1\le j\le r\).
In ordinary notation \(f_j=f^{(j)}\).  Define the cellwise quantities
\begin{equation}
 J_j=\sum_{k=a}^{b-1}\int_0^1P_j(t)f_j(k+t)\dd t,
 \qquad 0\le j\le r.\label{eq:Jdef}
\end{equation}
Every integral exists by the explicit product modulus and the
construction in Section~\ref{sec:integral}.  Translation of a rational
variable preserves the given moduli.  Finite additivity gives
\(J_0=\int_a^b f_0\).

Define the cellwise trapezoidal sum
\begin{equation}
 T=\frac12\sum_{k=a}^{b-1}\bigl(f_0(k)+f_0(k+1)\bigr).\label{eq:Tdef}
\end{equation}
For \(b>a\), this is
\(\tfrac12f_0(a)+\sum_{a<k<b}f_0(k)+\tfrac12f_0(b)\).
The cellwise definition makes \(T=0\) when \(a=b\).

Integration by parts on \([0,1]\), using \(P_j'=P_{j-1}\), gives
\begin{align}
 &\int_0^1P_j(t)f_j(k+t)\dd t
  +\int_0^1P_{j-1}(t)f_{j-1}(k+t)\dd t\notag\\
 &\hspace{2cm}=P_j(1)f_{j-1}(k+1)-P_j(0)f_{j-1}(k).\label{eq:cellparts}
\end{align}
For \(j=1\), the right side is
\((f_0(k+1)+f_0(k))/2\), so
\begin{equation}
 T=J_0+J_1.\label{eq:EMbase}
\end{equation}
For \(j\ge2\), the endpoints of \(P_j\) agree; summing the right
side of \eqref{eq:cellparts} telescopes.  Thus
\begin{equation}
 J_{j-1}+J_j=c_j\bigl(f_{j-1}(b)-f_{j-1}(a)\bigr).
 \label{eq:EMrecurrence}
\end{equation}

Substitute these identities successively into \eqref{eq:EMbase}.
Equivalently, add \eqref{eq:EMrecurrence} with alternating signs.
The intermediate \(J_j\)'s cancel and give
\begin{equation}
 T=J_0+\sum_{j=2}^r(-1)^jc_j
       \bigl(f_{j-1}(b)-f_{j-1}(a)\bigr)+(-1)^{r+1}J_r.
 \label{eq:EMgeneral}
\end{equation}
For \(r=1\), the middle sum is empty and this is exactly
\eqref{eq:EMbase}.  For \(r=2m\), all odd \(c_j\), \(j\ge3\),
vanish.  Adding \((f_0(a)+f_0(b))/2\) to \(T\) gives the inclusive
sum, proving \eqref{eq:emintro}--\eqref{eq:rintro}.

When \(a=b\), all cellwise sums and integrals, and all derivative
differences, are zero.  The endpoint term equals the one-term inclusive
sum, so the formula remains valid.  Negative integer endpoints cause
no change; signed integer arithmetic and the nonnegative cell count
\(b-a\) are primitive recursive.

The notation \(P_r(\{x\})\) can be introduced after the proof as an
abbreviation for the cellwise integrands.  We do not compute the floor
of a real name.  In particular, the jump of \(P_1(\{x\})\) at an
integer causes no continuity assumption across cells.  At orders
\(r\ge2\), the polynomial endpoints match anyway.

Section~\ref{sec:finiteEM} now supplies the all-orders rational proof
behind this convenient analytic notation.  Its purpose is to ensure
that \eqref{eq:EMgeneral} is never justified by applying
quantifier-free induction to an equality whose meaning contains an
unbounded precision quantifier.

\section{A fully rational Euler--Maclaurin certificate}\label{sec:finiteEM}

Assume \(L=b-a\ge1\) and \(r\ge1\).  Let an integer \(M\ge1\)
bound all \(f_j\), \(0\le j\le r\); compute it by a finite maximum
of the supplied bounds.  Fix a local precision \(q\) and choose
\begin{equation}
 \ell\ge\max\{q,\nu_0(q),\ldots,\nu_{r-1}(q)\},
 \quad H=2^\ell,\quad h=2^{-\ell},\quad p=q+\ell+4.
 \label{eq:EMschedulelocal}
\end{equation}
Here \(\nu_{j-1}\) is the modulus for \(f_{j-1}'=f_j\).
Put \(\eps=2^{-q}\) and use one common sample precision throughout:
\(\widehat f_j(x)=F_j(x,p)\).  In particular,
\(2^{-p}=\eps h/16\).  Formula (D) gives
\begin{align}
 \abs{\Delta\widehat f_{j-1}-h\widehat f_j}
 &\le h\eps+(2+h)2^{-p}\le\tfrac{19}{16}h\eps,
 \label{eq:rationalderiv}\\
 \abs{\widehat f_j}&\le M+1,&
 \abs{\Delta\widehat f_{j-1}}&\le h(M+3).
 \label{eq:rationalincrement}
\end{align}
All evaluations in these statements are at rational points
\(x=k+ih\), with \(i<H\).

Set
\begin{equation}
 C=2^r,\quad D_1=r2^r,\quad D_2=r^22^r,\quad
 D=2^r(M+1)(r+1)^2.\label{eq:Dconstants}
\end{equation}
Equations \eqref{eq:coefbound} and \eqref{eq:monomialremainder}
prove, for every \(j\le r\) and \(0\le t\le t+h\le1\),
\begin{equation}
 \abs{P_j(t)}\le C,\qquad
 \abs{\Delta P_j}\le D_1h,\qquad
 \abs{\Delta P_j-hP_{j-1}(t)}\le D_2h^2.
 \label{eq:Pfinitebounds}
\end{equation}
For the last expression take \(j\ge1\).

Apply \eqref{eq:productfinite} to
\(P_j(t)\widehat f_{j-1}(k+t)\), and subtract its proposed first-order
term.  The exact rational defect on one fine cell is
\begin{align}
 d_{j,k,i}={}&\Delta(P_j\widehat f_{j-1})
  -h(P_j\widehat f_j+P_{j-1}\widehat f_{j-1})\notag\\
 ={}&P_j(\Delta\widehat f_{j-1}-h\widehat f_j)
 +\widehat f_{j-1}(\Delta P_j-hP_{j-1})
 +\Delta P_j\Delta\widehat f_{j-1}.\label{eq:localdefect}
\end{align}
Use \(h\le\eps\) from \eqref{eq:EMschedulelocal}.  Its absolute value
is at most
\[
 h\eps\left(\tfrac{19}{16}C+(M+1)D_2+(M+3)D_1\right)
 \le hD\eps.
\]
For transparency, the last comparison follows from the rational identity
\[
 (M+1)(r+1)^2-
 \left(\tfrac{19}{16}+(M+1)r^2+(M+3)r\right)
 =r(M-1)+M-\tfrac3{16}>0.
\]

Define the rational quadratures
\begin{equation}
 Q_j=h\sum_{k=a}^{b-1}\sum_{i=0}^{H-1}
           P_j(ih)\widehat f_j(k+ih).\label{eq:Q}
\end{equation}
Let \(b_1\) be the cellwise trapezoidal sum of \(\widehat f_0\), and
for \(2\le j\le r\) put
\[
 b_j=c_j\bigl(\widehat f_{j-1}(b)-\widehat f_{j-1}(a)\bigr).
\]
For every \(1\le j\le r\), define
\(E_j=b_j-Q_{j-1}-Q_j\).
Summing \eqref{eq:localdefect} over \(H\) fine cells in each of \(L\)
unit cells telescopes exactly, because repeated endpoints use the same
sample.  Therefore
\begin{equation}
 \abs{E_j}\le LD\eps\quad(1\le j\le r).\label{eq:Ejbound}
\end{equation}
The following is an exact identity of rational expressions:
\begin{equation}
 b_1=Q_0+\sum_{j=2}^r(-1)^j b_j+(-1)^{r+1}Q_r
               +\sum_{j=1}^r(-1)^{j+1}E_j.\label{eq:finiteEMidentity}
\end{equation}
It follows directly by substituting the definition of \(E_j\);
each interior \(Q_j\) has coefficient zero.  A finite accumulator
proof establishes the same identity by quantifier-free induction.
Combining with \eqref{eq:Ejbound} gives the literal rational estimate
\begin{equation}
 \abs{b_1-Q_0-\sum_{j=2}^r(-1)^jb_j-(-1)^{r+1}Q_r}
 \le LrD2^{-q}.\label{eq:finiteEMbound}
\end{equation}
This is the central arithmetic theorem.  Its proof uses no real names,
integral objects, or quantified-function induction.

\subsection{From the finite certificate to the represented formula}

Also require \(\ell\ge\mu_0(q),\mu_r(q)\).  Then
\eqref{eq:quaderror}, applied cellwise, gives
\[
 \abs{Q_0-J_0}\le L(\eps+2^{-p})\le2L\eps.
\]
For the last product, whenever \(0\le t-t_i\le h\),
\[
 \abs{P_r(t)f_r(k+t)-P_r(t_i)f_r(k+t_i)}
 \le C\eps+MD_1h.
\]
Its quadrature error, including sample errors, is consequently at most
\[
 L\bigl((C+MD_1)\eps+C2^{-p}\bigr)
 \le LD\eps.
\]
Here \(D\ge2C+MD_1\) follows directly from \eqref{eq:Dconstants}.
The trapezoidal samples have total weight \(L\), so replacing them by
their names costs at most \(L2^{-p}\).  The endpoint corrections cost
at most \(2(r-1)C2^{-p}\).  Thus the error of
\eqref{eq:EMgeneral}, as a represented identity, is bounded by
\begin{align*}
 \bigl(LrD+2L+LD+L+2(r-1)C\bigr)\eps
 &\le7LrD\eps\\
 &<8LrD\eps.
\end{align*}
These comparisons can be stated at fixed indices of the output names,
adding their radius terms, and then justified by the rational closure
lemma.  No existence of limits is being invoked.

For a requested output precision \(n\), one completely explicit
schedule is
\begin{align}
 q&=n+\sigma\bigl(8(L+1)(r+1)D\bigr),\notag\\
 \ell&=\max\{q,\mu_0(q),\mu_r(q),
                         \nu_0(q),\ldots,\nu_{r-1}(q)\},\notag\\
 p&=q+\ell+4.\label{eq:EMfinalschedule}
\end{align}
Then \(8LrD2^{-q}<2^{-n}\).  The maximum is a bounded scan over the
given derivative family.  All coefficient lists, evaluation indices,
and grid sizes are primitive recursive.  The order \(r\) can be a
number variable throughout this proof.  Only after the finite rational
argument is complete do we universally generalize the requested
precision \(n\).  This resolves the quantifier-free-induction issue
for the all-orders Euler--Maclaurin theorem.

\section{An elementary bound on the remainder}\label{sec:remainder}

A convenient bound using only rational arithmetic is
\begin{equation}
 \abs{P_n(t)}\le2^{-n}\qquad(0\le t\le1).\label{eq:halfbound}
\end{equation}
It is not claimed to be sharp.  It avoids the \(\pi\), zeta values,
and Fourier estimates that occur in a customary sharper bound.

There is a short explanation of the contraction.  The normalization
in \eqref{eq:Pdef} gives the polynomial identity
\[
 P_n(t)=\int_0^t uP_{n-1}(u)\dd u
            -\int_t^1(1-u)P_{n-1}(u)\dd u.
\]
It follows either by coefficientwise calculation or by differentiating
both sides and comparing their zero means.  A bound \(K\) on
\(P_{n-1}\) therefore gives
\[
 \abs{P_n(t)}\le\tfrac K2\bigl(t^2+(1-t)^2\bigr)\le K/2.
\]
To justify uniform induction on \(n\) strictly within the stated
foundation, we give a finite coefficient certificate for this argument.

\subsection{The Bernstein coefficient certificate}

The degree-\(d\) Bernstein basis is
\[
 \beta_{d,i}(t)=\binom di t^i(1-t)^{d-i}\quad(0\le i\le d).
\]
For rational \(t\in[0,1]\), all basis values are nonnegative and
their sum is one, by the finite binomial identity.  Therefore a finite
coefficient bound \(\abs{b_i}\le K\) implies
\(\abs{\sum_i b_i\beta_{d,i}(t)}\le K\).

Suppose \(P_{n-1}=\sum_{i=0}^{n-1}b_i\beta_{n-1,i}\).  The degree-\(n\)
Bernstein coefficients of its normalized antiderivative \(P_n\) are
\begin{equation}
 d_j=\frac{\displaystyle\sum_{i<j}(i+1)b_i
                   -\displaystyle\sum_{i\ge j}(n-i)b_i}
                  {n(n+1)},\qquad 0\le j\le n.\label{eq:Bernsteinrec}
\end{equation}
To verify this, note that an antiderivative has coefficients
\(a_j=n^{-1}\sum_{i<j}b_i\): differentiating the Bernstein basis
gives \(n(a_{i+1}-a_i)=b_i\).  The coefficientwise mean of a
degree-\(n\) Bernstein polynomial is the average of its \(n+1\)
coefficients.  To prove the basis mean formula entirely by finite
polynomial identities, for \(0\le j<n\) differentiate to obtain
\[
 \beta'_{n+1,j+1}=(n+1)(\beta_{n,j}-\beta_{n,j+1}).
\]
The polynomial on the left before differentiation vanishes at both
endpoints, so coefficientwise integration shows that adjacent basis
means are equal.  Their sum is the integral of
\(\sum_j\beta_{n,j}=1\), namely one.  Hence every basis mean is
\(1/(n+1)\).  For \(n=0\), the statement is immediate.
Subtracting this mean from each \(a_j\) gives
\eqref{eq:Bernsteinrec}.

The sum of the absolute weights in \eqref{eq:Bernsteinrec} is
\begin{equation}
 \frac{j(j+1)+(n-j)(n-j+1)}{2n(n+1)}\le\frac12,\label{eq:Bernsteinweights}
\end{equation}
because the numerator is \(n(n+1)-2j(n-j)\).
Starting with the single coefficient \(1\) of \(P_0\), the maximum
absolute Bernstein coefficient thus contracts by \(1/2\) at every
step.  The inductive invariant is the bounded finite-list test
\(\max_i\abs{b_i}\le2^{-(n-1)}\), expressed as a rational inequality;
the coefficient identity is another finite-list equality.  Both are
quantifier-free after coding.  Nonnegativity and partition of unity
then prove \eqref{eq:halfbound} for any rational \(t\).
The name extension proves the same represented inequality on the
whole interval, without further induction.

\begin{corollary}[Explicit remainder estimate]\label{cor:remainder}
If \(\abs{f_{2m}}\le M_{2m}\) on \([a,b]\), then
\begin{equation}
 \abs{R_m}\le4^{-m}\int_a^b\abs{f_{2m}(x)}\dd x
              \le(b-a)4^{-m}M_{2m}.\label{eq:remainderbound}
\end{equation}
\end{corollary}
\begin{proof}
The absolute-value function preserves the continuity modulus, by
\(\bigl|\abs u-\abs v\bigr|\le\abs{u-v}\), so its integral is
available.  Apply \eqref{eq:halfbound} on each unit cell, the integral
triangle inequality (the finite-sum proof of
Proposition~\ref{prop:integralrules}), and finite additivity.
\end{proof}

This estimate is a bound at a specified order.  Since the input bounds
\(M_{2m}\) can grow faster than \(4^m\), it implies no general
convergence of an infinite Euler--Maclaurin series.  For a polynomial
of degree less than \(2m\), the top derivative is zero and the
remainder vanishes exactly.

\section{Rational meshes and an exact example}\label{sec:examples}

Let \(h>0\) be rational, \(N\in\N\), and \(b=a+Nh\), with rational
\(a\).  Apply the unit-mesh formula to \(g(t)=f(a+ht)\), whose
certified derivatives are \(g_j(t)=h^jf_j(a+ht)\).  The derivative
certificates follow by direct substitution in (D) and explicit
precision shifts for the rational factors.  Finite sums also prove
the integral change of variable.  Multiplying the unit-mesh formula
by \(h\) gives
\begin{align}
 h\sum_{k=0}^Nf(a+kh)
 ={}&\int_a^b f(x)\dd x+\frac h2(f(a)+f(b))\notag\\
 &+\sum_{s=1}^m h^{2s}c_{2s}
          \bigl(f_{2s-1}(b)-f_{2s-1}(a)\bigr)+R_{m,h},
 \label{eq:scaledEM}\\
 R_{m,h}={}&-h^{2m+1}\sum_{k=0}^{N-1}\int_0^1
       P_{2m}(t)f_{2m}(a+h(k+t))\dd t.\label{eq:scaledR}
\end{align}
The remainder obeys
\[
 \abs{R_{m,h}}\le (b-a)h^{2m}4^{-m}M_{2m}.
\]
For \(N=0\), the identity is interpreted exactly as in the zero-length
unit-mesh case.

A useful sign check has a \emph{nonzero} remainder.  Let
\(f(x)=x^4\), \(a=0\), \(b=3\), and \(m=1\).  Direct polynomial
integration gives
\[
 \int_0^3x^4\dd x=\frac{243}{5},\qquad
 \frac{f(0)+f(3)}2=\frac{81}{2},\qquad
 c_2(f'(3)-f'(0))=9.
\]
For the remainder, \(f''(x)=12x^2\) and
\[
 \int_0^1P_2(t)\dd t=0,\quad
 \int_0^1tP_2(t)\dd t=0,\quad
 \int_0^1t^2P_2(t)\dd t=\frac1{360}.
\]
Each of the three cells therefore contributes \(1/30\) before the
minus sign in \eqref{eq:rintro}, so \(R_1=-1/10\).  The exact identity
is
\[
 \boxed{\displaystyle
 \sum_{k=0}^{3}k^4=98=\frac{243}{5}+\frac{81}{2}+9-\frac1{10}.}
\]
This checks the endpoint convention and the negative sign of the even
remainder without relying on exact truncation.

\section{Why the proofs fit \texorpdfstring{\ipra}{iPRA}}\label{sec:audit}

The proof-theoretic claim is a human-readable formalizability argument,
with the following concrete accounting.  No analytic lemma is being
assumed as an axiom of \ipra.

\begin{longtable}{>{\raggedright\arraybackslash}p{.26\textwidth}>{\raggedright\arraybackslash}p{.64\textwidth}}
\toprule
Construction or proof & Primitive recursive operation and permitted induction\\\midrule
\endhead
Rational arithmetic & Signed-integer operations, cross multiplication,
and positive-denominator coding.\\
Precision choices & Powers of two, bounded search \(\sigma\), and finite
maxima of the supplied moduli.\\
Integral evaluator & A finite grid followed by a finite rational sum;
the grid length is the explicit power of two in
\eqref{eq:integralname}.\\
Integral coherence & Finite refinement or overlap weights, positivity,
and the finite triangle inequality.\\
FTC I & The fixed-index rational estimates
\eqref{eq:addfinite} and \eqref{eq:localintegral}.\\
FTC II & The partial-sum invariant behind
\eqref{eq:finiteFTC}; internal precision includes its \(2N\) factor.\\
Riemann FTC II & The supplied finite partition is checked by bounded
tests; the same partial-sum argument gives the \(4N\) factor.\\
Bernoulli identities & Equality of finite coefficient lists, with
coefficientwise derivative, reflection, and mean operations.\\
Euler--Maclaurin & The exact rational invariant
\eqref{eq:finiteEMidentity} and its error bound
\eqref{eq:finiteEMbound}, with precision fixed before induction.\\
Remainder bound & A finite Bernstein coefficient maximum and the
explicit contraction \eqref{eq:Bernsteinweights}.\\
Name equalities & Universal generalization over number precisions after
the finite proof, plus the rational closure lemma when eliminating
auxiliary errors.\\\bottomrule
\end{longtable}

For example, the induction invariant for a finite triangle estimate
may be written
\[
 \operatorname{valid}(z,w,n)=1\ \longrightarrow\
 \abs{\sum_{i<j}z_i}\le\sum_{i<j}w_i\quad(j\le n),
\]
where \(\operatorname{valid}\) is a primitive recursive Boolean
checking all \(\abs{z_i}\le w_i\).  After coding, this is a
quantifier-free formula with the array codes and \(n\) as parameters.
The successor step is the triangle inequality; no quantified analytic
assertion occurs in the induction formula.

The same principle applies to the coefficient tables.  It is
insufficient to say only that a proof is finite: every formal proof
is finite.  What matters is that the inductions displayed here concern
rational equalities, rational inequalities, and bounded finite-list
tests.  The all-orders construction in Section~\ref{sec:finiteEM}
specifically prevents the tempting invalid step of inducting on
``the two real names are equal at every precision.''

The inputs' regularity promises are hypotheses of the schema.
For a concrete elementary function, those promises must themselves be
proved from its primitive recursive definition before the result is
unconditional.  For rational polynomials they follow from
\eqref{eq:monomialremainder} and finite coefficient bounds.  The
proofs do not provide a general procedure deciding whether an arbitrary
program is a valid continuous or differentiable function name.

\section{Lean formalization: existing foundations and completion path}\label{sec:lean}

\subsection{Pinned sources and the trust boundary}

The inspected ProveIt revision is
\begin{center}
\code{0bf379a84212abd05e9a3f0b1320a3d7da509cf9}.
\end{center}
Its root toolchain is Lean \code{v4.32.0}, with Mathlib
\code{v4.32.0}.  The useful deeply embedded foundations are the
\code{PR}, \code{Deduction}, \code{Syntax}, \code{BiInterpretation},
and \code{NativeBiInterpretation} modules
\cite{prsource,nativesource}.  Their source directory is
\begin{quote}\small
\path{Logic/PrimitiveRecursiveArithmetic/ConciseIPRA/Lean/ConciseIPRA/}
\end{quote}

In particular, \code{PR} is an arity-indexed syntax generated by zero,
successor, projections, composition, and primitive recursion.
\code{PRAAx} contains defining equations and quantifier-free induction;
\code{PRA.qfInd} constructs the corresponding object-theory induction
proof.  The final formal goal must be an actual derivability statement
\[
 \Gamma\vdash[\code{PRAAx}]\varphi,
\]
in namespace \code{LeanProofs.ConciseIPRA}.  Here \(\Gamma\) is the
arithmetized list of input promises and \(\varphi\) is a coded
precision inequality or its universal closure.  Proving a proposition
in ambient Lean about rational or real functions is a useful prototype,
but it is not a proof of this derivability statement.

The repository also contains an existing first-order sawtooth
Euler--Maclaurin theorem in its \code{IntegerPoints} development
\cite{existingem}:
\begin{quote}\small
\code{LeanProofs.IntegerPoints.EM.sum\_eq\_integral}.
\end{quote}
It provides a useful independent check of signs and
endpoint conventions.  It uses real floor, interval integrals, measure
theory, and compactness to obtain a bound.  It is not an \ipra\ proof
dependency.  Likewise, Mathlib's standard FTC development is an
ordinary real-analysis development \cite{mathlibftc}; importing it
would not settle the foundational claim here.

The listing of \path{docs/drafts} at this revision included FTA, Basel,
isoperimetric, Minkowski, prime-number-theorem, and Stirling archives.
Their binary contents were not audited and no result from them is used
as a dependency.  The mathematical proof in this package is independent
of those drafts.

\subsection{Delivered Lean candidate}

The file \path{lean/FiniteSkeleton.lean} contains candidate proof scripts
for rational finite-sum linearity, telescoping, the triangle inequality,
the finite FTC bound, and an alternating finite recurrence.  For example,
its central FTC statement is the following arithmetic implication:
\[
 \bigl(\forall i<n,\ \abs{F_{i+1}-F_i-hf_i}\le e\bigr)
 \ \Longrightarrow\
 \abs{F_n-F_0-h\sum_{i<n}f_i}\le ne.
\]
It also includes structures specifying regular names and uniform
rational approximations.  Those structures are metatheoretic interfaces;
they do not assert that every Lean function in them is primitive
recursive.

\textbf{Status.} The candidate contains no \code{sorry} declarations or
new axioms, but it is \emph{uncompiled}.  No Lean, Lake, or Elan
executable was available in the preparation environment.  The file
therefore makes no claim of successful elaboration or kernel checking.
The full analytic theorems and their deeply embedded \ipra\ proofs
remain to be implemented.  The exact Python calculations are validation
artifacts, not substitutes for that implementation.

\subsection{A concrete path to a complete formalization}

The accompanying \path{lean/ROADMAP.md} expands the following sequence.

\begin{enumerate}
\item \textbf{Compile the finite semantic prototype.}
Use the pinned checkout and its dependencies, and run
\code{lake env lean} on \path{FiniteSkeleton.lean}.  Resolve any
elaboration issues and inspect the declarations' axioms.  Extend it
with the exact finite identity \eqref{eq:finiteEMidentity} and its
rational estimate.  At this stage there are still no object-theory
proofs.

\item \textbf{Build the encoded arithmetic layer.}
Give fixed \code{PR} expressions for signed rationals, lists,
polynomial operations, powers of two, sums, bounded maxima, and
\(\sigma\).  Prove their identities by constructing
\code{Prf PRAAx} derivations.  Use proof-producing macros for
rational ring equalities, cross-multiplied inequalities, and finite
triangle inequalities.  Each macro must return a proof in the
embedded calculus, rather than only an ambient Lean proof of truth.

\item \textbf{Deeply embed the data promises.}
Write formulas for (C), (B), (U), and (D), using the chosen fixed
evaluator and modulus symbols.  Encode interval membership and finite
partition validity by decidable tests.  Keep universal data assumptions
in the sequent context.  Replace bounded quantifiers inside induction
invariants by Boolean folds.

\item \textbf{Formalize rational integration and FTC.}
Implement \eqref{eq:integralname}.  Derive
\eqref{eq:refine}, \eqref{eq:integralcauchy},
\eqref{eq:addfinite}, the literal derivative certificate in
Theorem~\ref{thm:ftc1}, and \eqref{eq:finiteFTC}.
The schedules \eqref{eq:FTCprecision} produce the final inequalities.
The optional general version adds only the finite tagged-partition
interface and \eqref{eq:taggedfinite}.

\item \textbf{Formalize the Bernoulli table.}
Use \eqref{eq:Pcoeff} as the recursion on a finite rational coefficient
list.  Prove the derivative, reflection, endpoint, and odd-vanishing
identities as list equalities.  Prove the Bernstein recursion
\eqref{eq:Bernsteinrec} and its finite maximum bound.  No formal power
series or classical real integration is necessary.

\item \textbf{Prove the uniform finite EM certificate.}
Implement the parameters \eqref{eq:Dconstants} and
\eqref{eq:EMfinalschedule}.  Prove
\eqref{eq:localdefect}, \eqref{eq:Ejbound},
\eqref{eq:finiteEMidentity}, and \eqref{eq:finiteEMbound}
by quantifier-free finite accumulator invariants.  Only afterward
universally generalize output precision and identify the named
integrals.  This stage also gives the general odd-order version.

\item \textbf{Transfer presentations if desired.}
After the \code{Prf PRAAx} terms exist, apply the already defined
\code{pra\_iff\_native} from
\code{NativeBiInterpretation.lean} to obtain native scalar-core
proofs.  That theorem preserves the explicit hypotheses; it does
not extract \ipra\ proofs automatically from ambient Lean analysis.
\end{enumerate}

The missing work is substantial formal encoding and proof production,
but the mathematical dependencies, exact intermediate statements,
permitted induction invariants, and precision algorithms are all
specified above.  There is no remaining appeal to a qualitative
analytic theorem that would need to be smuggled into \ipra.

\section{Artifacts and reproducibility}\label{sec:artifacts}

The archive contains this \LaTeX\ source and its compiled PDF, the
Lean candidate and roadmap, exact Python implementations, a validation
record, and source provenance.  Rebuild instructions are in
\path{README.md}.

The standard-library-only program \path{code/bernoulli_em.py} uses
\code{fractions.Fraction} throughout.  It computes normalized Bernoulli
polynomials, their Bernstein coefficients, polynomial derivatives and
integrals, exact Euler--Maclaurin terms and remainders, and finite-grid
product defects with rounded rational oracles.  The recorded run passed:
\begin{itemize}
\item 382 Bernoulli identities and coefficient-certificate checks
through order 20;
\item 85 polynomial FTC checks;
\item 612 unit-mesh Euler--Maclaurin identities, including 100 cases
with nonzero remainder and negative, shifted, and zero-length intervals;
\item 272 rational-mesh identities;
\item 600 local finite-product checks in 16 rounded-oracle cases,
using the sample precision \(p=q+\ell+4\).
\end{itemize}
The companion \path{code/integral_demo.py} exercises the canonical
integrator on the rationally evaluable nonpolynomial function
\(f(x)=1/(1+x)\) on \([0,1]\), and checks finite additivity and
derivative-certificate instances with exact arithmetic.
These are regression and implementation checks.  The proofs of the
universal statements are the arguments in the preceding sections;
formal \ipra\ certification is the explicit future Lean task just
described.

\begin{thebibliography}{9}
\small\raggedright
\bibitem{prsource}
Vladimir Reshetnikov, \emph{ProveIt: ConciseIPRA, PR.lean},
revision \code{0bf379a84212abd05e9a3f0b1320a3d7da509cf9}.
\url{https://github.com/VladimirReshetnikov/ProveIt/blob/0bf379a84212abd05e9a3f0b1320a3d7da509cf9/Logic/PrimitiveRecursiveArithmetic/ConciseIPRA/Lean/ConciseIPRA/PR.lean}.

\bibitem{nativesource}
Vladimir Reshetnikov, \emph{ProveIt: ConciseIPRA,
NativeBiInterpretation.lean}, same revision.
\url{https://github.com/VladimirReshetnikov/ProveIt/blob/0bf379a84212abd05e9a3f0b1320a3d7da509cf9/Logic/PrimitiveRecursiveArithmetic/ConciseIPRA/Lean/ConciseIPRA/NativeBiInterpretation.lean}.

\bibitem{existingem}
Vladimir Reshetnikov, \emph{ProveIt: IntegerPoints,
EulerMaclaurin.lean}, same revision.
\url{https://github.com/VladimirReshetnikov/ProveIt/blob/0bf379a84212abd05e9a3f0b1320a3d7da509cf9/NumberTheory/IntegerPoints/Lean/IntegerPoints/EulerMaclaurin.lean}.

\bibitem{dlmfem}
NIST Digital Library of Mathematical Functions,
\S2.10(i), \emph{Euler--Maclaurin Formula}.
\url{https://dlmf.nist.gov/2.10.i}.

\bibitem{dlmfbernoulli}
NIST Digital Library of Mathematical Functions,
\S24.2, \emph{Definitions and Generating Functions}.
\url{https://dlmf.nist.gov/24.2}.

\bibitem{dlmfproperties}
NIST Digital Library of Mathematical Functions,
\S24.4, \emph{Basic Properties}.
\url{https://dlmf.nist.gov/24.4}.

\bibitem{mathlibftc}
Mathlib documentation, \emph{Fundamental Theorem of Calculus}.
\url{https://leanprover-community.github.io/mathlib4_docs/Mathlib/MeasureTheory/Integral/IntervalIntegral/FundThmCalculus.html}.
\end{thebibliography}
\end{document}
